% Link-Cut Tree

A data structure that maintains a forest of rooted trees under dynamic edge insertions (\lstinline|link|) and deletions (\lstinline|cut|). All operations run in $O(\log N)$ amortized time using Splay Trees.

\begin{lstlisting}
struct Node {
    int p = 0, ch[2] = {0, 0};
    bool rev = false;
};

vector<Node> t;

void init_lct(int n) {
    t.assign(n + 1, Node());
}

bool is_root(int x) {
    return t[t[x].p].ch[0] != x && t[t[x].p].ch[1] != x;
}

void push(int x) {
    if (x && t[x].rev) {
        swap(t[x].ch[0], t[x].ch[1]);
        if (t[x].ch[0]) t[t[x].ch[0]].rev ^= 1;
        if (t[x].ch[1]) t[t[x].ch[1]].rev ^= 1;
        t[x].rev = false;
    }
}

void push_all(int x) {
    if (!is_root(x)) push_all(t[x].p);
    push(x);
}

void rotate(int x) {
    int y = t[x].p, z = t[y].p;
    int k = (t[y].ch[1] == x);
    if (!is_root(y)) t[z].ch[t[z].ch[1] == y] = x;
    t[x].p = z;
    t[y].ch[k] = t[x].ch[k ^ 1];
    if (t[x].ch[k ^ 1]) t[t[x].ch[k ^ 1]].p = y;
    t[x].ch[k ^ 1] = y;
    t[y].p = x;
}

void splay(int x) {
    push_all(x);
    while (!is_root(x)) {
        int y = t[x].p, z = t[y].p;
        if (!is_root(y)) {
            if ((t[y].ch[1] == x) ^ (t[z].ch[1] == y)) rotate(x);
            else rotate(y);
        }
        rotate(x);
    }
}

void access(int x) {
    for (int t_ch = 0, r = x; r; t_ch = r, r = t[r].p) {
        splay(r);
        t[r].ch[1] = t_ch;
    }
    splay(x);
}

void make_root(int x) {
    access(x);
    t[x].rev ^= 1;
}

int find_root(int x) {
    access(x);
    while (t[x].ch[0]) {
        push(x);
        x = t[x].ch[0];
    }
    splay(x);
    return x;
}

void link(int x, int y) {
    make_root(x);
    if (find_root(y) != x) t[x].p = y;
}

void cut(int x, int y) {
    make_root(x);
    if (find_root(y) == x && t[y].p == x && !t[y].ch[0]) {
        t[y].p = t[x].ch[1] = 0;
    }
}
\end{lstlisting}

\begin{itemize}
    \item Call \lstinline|init_lct(n)| to allocate memory for $n$ nodes ($1$-indexed).
    \item Call \lstinline|link(x, y)| to add an undirected edge between $x$ and $y$. It automatically checks connectivity to prevent cycles.
    \item Call \lstinline|cut(x, y)| to remove the edge between $x$ and $y$ if it exists.
    \item Call \lstinline|make_root(x)| to change the root of the tree containing $x$ to be $x$.
    \item Call \lstinline|find_root(x)| to get the actual component root of $x$. Useful for dynamic connectivity checks via \lstinline|find_root(x) == find_root(y)|.
\end{itemize}
